\documentclass[12pt, ukrainian]{article}
\usepackage[a4paper]{geometry}
\setlength\parindent{0mm}
%\setlength\parskip{4mm}
\geometry{top=1.3cm,bottom=1.3cm,left=1.5cm,right=1.5cm}
\pagestyle{empty}
\usepackage[T2A]{fontenc}
\usepackage[utf8]{inputenc}
\usepackage[ukrainian]{babel}
\begin{document}
%begin
{{\begin {scriptsize}\par
{ \center  Київський національний університет імені Тараса Шевченка \par}
{\parbox {5 cm}{Освітня програма \hfill}{\parbox {7 cm}{Прикладна математика\hfill}}\parbox {6 cm}{\hfill Семестр 1}}\par
{\parbox {5 cm} {Навчальний предмет \hfill}\parbox {7 cm}{Програмування \hfill}\parbox {6cm}{\hfill}\par}\vspace{-15pt} \end{scriptsize}}{\center Екзаменаційний білет \No \textbf { 1 }\par}}
{%beginitem
1. Вбудовані числові типи мови Python.%enditem
\par
%beginitem
2. Написати програму, що запитує у користувача значення дійсних параметрів $x$ та $y$, обчислює для них значення виразу 
$$\frac{\log_{2} x + 60}{(\arccos x - 0.3) (y - 11)}$$ 
та повiдомляє введенi данi та результат обчислення.
 Оцінюється правильність та якість коду.

%enditem
\par
%beginitem
3. Написати функцію, що для інтервалу $[a, b)$ знаходить знаходить число з найменшим добутком простих дільників. Якщо таких чисел кілька, повертає найбільше.  Використання рядків та структур даних заборонено. Якість та швидкодія коду також оцінюються. 

%enditem
}
{\smallskip\smallskip \begin {scriptsize} Затверджено на засіданні кафедри теоретичної кібернетики, протокол \No 1  від   06.09.2023 р.\par\smallskip\smallskip\smallskip\smallskip
{\parbox {6.5 cm} {Зав. кафедри \dotfill Ю.В.~Крак}}\hspace {0.7cm} {\hbox to 6.5 cm { Екзаменатор \dotfill Т.О.~Карнаух}} \end{scriptsize}}
%end
\vspace{0.9cm}
%begin
{{\begin {scriptsize}\par
{ \center  Київський національний університет імені Тараса Шевченка \par}
{\parbox {5 cm}{Освітня програма \hfill}{\parbox {7 cm}{Прикладна математика\hfill}}\parbox {6 cm}{\hfill Семестр 1}}\par
{\parbox {5 cm} {Навчальний предмет \hfill}\parbox {7 cm}{Програмування \hfill}\parbox {6cm}{\hfill}\par}\vspace{-15pt} \end{scriptsize}}{\center Екзаменаційний білет \No \textbf { 2 }\par}}
{%beginitem
1. Цикл for, вигляд і семантика інструкції циклу for.%enditem
\par
%beginitem
2. Написати програму, що запитує у користувача значення дійсних параметрів $x$ та $y$, обчислює для них значення виразу 
$$\frac{\ln x - 82}{(\sin x + 0.4) (y + 89)}$$ 
та повiдомляє введенi данi та результат обчислення.
 Оцінюється правильність та якість коду.

%enditem
\par
%beginitem
3. Написати функцію, що для інтервалу $[a, b)$ знаходить число з найбільшою кількістю різних простих дільників, що містять цифру 7 у своєму десятковому запису. Якщо їх кілька, то цікавить найбільше.  Використання рядків та структур даних заборонено. Якість та швидкодія коду також оцінюються. 

%enditem
}
{\smallskip\smallskip \begin {scriptsize} Затверджено на засіданні кафедри теоретичної кібернетики, протокол \No 1  від   06.09.2023 р.\par\smallskip\smallskip\smallskip\smallskip
{\parbox {6.5 cm} {Зав. кафедри \dotfill Ю.В.~Крак}}\hspace {0.7cm} {\hbox to 6.5 cm { Екзаменатор \dotfill Т.О.~Карнаух}} \end{scriptsize}}
%end
\vspace{0.9cm}
%begin
{{\begin {scriptsize}\par
{ \center  Київський національний університет імені Тараса Шевченка \par}
{\parbox {5 cm}{Освітня програма \hfill}{\parbox {7 cm}{Прикладна математика\hfill}}\parbox {6 cm}{\hfill Семестр 1}}\par
{\parbox {5 cm} {Навчальний предмет \hfill}\parbox {7 cm}{Програмування \hfill}\parbox {6cm}{\hfill}\par}\vspace{-15pt} \end{scriptsize}}{\center Екзаменаційний білет \No \textbf { 3 }\par}}
{%beginitem
1. Поняття типу даних.%enditem
\par
%beginitem
2. Написати програму, що запитує у користувача значення дійсних параметрів $x$ та $y$, обчислює для них значення виразу 
$$\frac{\log_{10} x + 45}{(\tg x - 0.8) (y + 44)}$$ 
та повiдомляє введенi данi та результат обчислення.
 Оцінюється правильність та якість коду.

%enditem
\par
%beginitem
3. Написати функцію, що для інтервалу $[a, b)$ знаходить число з найбільшою кількістю різних простих дільників. Якщо їх кілька, то цікавить найменше.  Використання рядків та структур даних заборонено. Якість та швидкодія коду також оцінюються. 

%enditem
}
{\smallskip\smallskip \begin {scriptsize} Затверджено на засіданні кафедри теоретичної кібернетики, протокол \No 1  від   06.09.2023 р.\par\smallskip\smallskip\smallskip\smallskip
{\parbox {6.5 cm} {Зав. кафедри \dotfill Ю.В.~Крак}}\hspace {0.7cm} {\hbox to 6.5 cm { Екзаменатор \dotfill Т.О.~Карнаух}} \end{scriptsize}}
%end
\newpage
%begin
{{\begin {scriptsize}\par
{ \center  Київський національний університет імені Тараса Шевченка \par}
{\parbox {5 cm}{Освітня програма \hfill}{\parbox {7 cm}{Прикладна математика\hfill}}\parbox {6 cm}{\hfill Семестр 1}}\par
{\parbox {5 cm} {Навчальний предмет \hfill}\parbox {7 cm}{Програмування \hfill}\parbox {6cm}{\hfill}\par}\vspace{-15pt} \end{scriptsize}}{\center Екзаменаційний білет \No \textbf { 4 }\par}}
{%beginitem
1. Семантика функції та її виклику, параметри та аргументи.%enditem
\par
%beginitem
2. Написати програму, що запитує у користувача значення дійсних параметрів $x$ та $y$, обчислює для них значення виразу 
$$\frac{\pi x - 67}{(\ctg x + 0.7) (y - 81)}$$ 
та повiдомляє введенi данi та результат обчислення.
 Оцінюється правильність та якість коду.

%enditem
\par
%beginitem
3. Написати функцію, що для інтервалу $[a, b)$ знаходить знаходить найбільший степінь числа 3, на який ділиться хоча б одне ціле число з цього інтервалу.  Використання рядків та структур даних заборонено. Якість та швидкодія коду також оцінюються. 

%enditem
}
{\smallskip\smallskip \begin {scriptsize} Затверджено на засіданні кафедри теоретичної кібернетики, протокол \No 1  від   06.09.2023 р.\par\smallskip\smallskip\smallskip\smallskip
{\parbox {6.5 cm} {Зав. кафедри \dotfill Ю.В.~Крак}}\hspace {0.7cm} {\hbox to 6.5 cm { Екзаменатор \dotfill Т.О.~Карнаух}} \end{scriptsize}}
%end
\vspace{0.9cm}
%begin
{{\begin {scriptsize}\par
{ \center  Київський національний університет імені Тараса Шевченка \par}
{\parbox {5 cm}{Освітня програма \hfill}{\parbox {7 cm}{Прикладна математика\hfill}}\parbox {6 cm}{\hfill Семестр 1}}\par
{\parbox {5 cm} {Навчальний предмет \hfill}\parbox {7 cm}{Програмування \hfill}\parbox {6cm}{\hfill}\par}\vspace{-15pt} \end{scriptsize}}{\center Екзаменаційний білет \No \textbf { 5 }\par}}
{%beginitem
1. Оператори порівняння (у широкому розумінні мови Python) та булеві оператори.%enditem
\par
%beginitem
2. Написати програму, що запитує у користувача значення дійсних параметрів $x$ та $y$, обчислює для них значення виразу 
$$\frac{\sqrt x + 78}{(\arctg x - 0.5) (y + 21)}$$ 
та повiдомляє введенi данi та результат обчислення.
 Оцінюється правильність та якість коду.

%enditem
\par
%beginitem
3. Написати функцію, яка для цілого числа знаходить кількість його різних простих дільників, що мають у десятковому запису цифру 1.  Використання рядків та структур даних заборонено. Якість та швидкодія коду також оцінюються. 

%enditem
}
{\smallskip\smallskip \begin {scriptsize} Затверджено на засіданні кафедри теоретичної кібернетики, протокол \No 1  від   06.09.2023 р.\par\smallskip\smallskip\smallskip\smallskip
{\parbox {6.5 cm} {Зав. кафедри \dotfill Ю.В.~Крак}}\hspace {0.7cm} {\hbox to 6.5 cm { Екзаменатор \dotfill Т.О.~Карнаух}} \end{scriptsize}}
%end
\vspace{0.9cm}
%begin
{{\begin {scriptsize}\par
{ \center  Київський національний університет імені Тараса Шевченка \par}
{\parbox {5 cm}{Освітня програма \hfill}{\parbox {7 cm}{Прикладна математика\hfill}}\parbox {6 cm}{\hfill Семестр 1}}\par
{\parbox {5 cm} {Навчальний предмет \hfill}\parbox {7 cm}{Програмування \hfill}\parbox {6cm}{\hfill}\par}\vspace{-15pt} \end{scriptsize}}{\center Екзаменаційний білет \No \textbf { 6 }\par}}
{%beginitem
1. Поняття часової та просторової складності обчислень. Складність у середньому. Асимптотичні оцінки складності.%enditem
\par
%beginitem
2. Написати програму, що запитує у користувача значення дійсних параметрів $x$ та $y$, обчислює для них значення виразу 
$$\frac{\ x - 92}{(\arcsin x + 0.1) (y - 62)}$$ 
та повiдомляє введенi данi та результат обчислення.
 Оцінюється правильність та якість коду.

%enditem
\par
%beginitem
3. Написати функцію, що для інтервалу $[a, b)$ знаходить найближчу пару чисел з цього інтервалу, що мають парну кількість різних простих дільників. Якщо таких пар кілька, то цікавить пара з найбільшою сумою.  Використання рядків та структур даних заборонено. Якість та швидкодія коду також оцінюються. 

%enditem
}
{\smallskip\smallskip \begin {scriptsize} Затверджено на засіданні кафедри теоретичної кібернетики, протокол \No 1  від   06.09.2023 р.\par\smallskip\smallskip\smallskip\smallskip
{\parbox {6.5 cm} {Зав. кафедри \dotfill Ю.В.~Крак}}\hspace {0.7cm} {\hbox to 6.5 cm { Екзаменатор \dotfill Т.О.~Карнаух}} \end{scriptsize}}
%end
\newpage
%begin
{{\begin {scriptsize}\par
{ \center  Київський національний університет імені Тараса Шевченка \par}
{\parbox {5 cm}{Освітня програма \hfill}{\parbox {7 cm}{Прикладна математика\hfill}}\parbox {6 cm}{\hfill Семестр 1}}\par
{\parbox {5 cm} {Навчальний предмет \hfill}\parbox {7 cm}{Програмування \hfill}\parbox {6cm}{\hfill}\par}\vspace{-15pt} \end{scriptsize}}{\center Екзаменаційний білет \No \textbf { 7 }\par}}
{%beginitem
1. Лексеми мови Python.%enditem
\par
%beginitem
2. Написати програму, що запитує у користувача значення дійсних параметрів $x$ та $y$, обчислює для них значення виразу 
$$\frac{\sqrt x - 23}{(\cos x - 0.6) (y - 32)}$$ 
та повiдомляє введенi данi та результат обчислення.
 Оцінюється правильність та якість коду.

%enditem
\par
%beginitem
3. Написати функцію, що для інтервалу $[a, b)$ знаходить знаходить число з найбільшим добутком простих дільників. Якщо таких чисел кілька, повертає найменше.  Використання рядків та структур даних заборонено. Якість та швидкодія коду також оцінюються. 

%enditem
}
{\smallskip\smallskip \begin {scriptsize} Затверджено на засіданні кафедри теоретичної кібернетики, протокол \No 1  від   06.09.2023 р.\par\smallskip\smallskip\smallskip\smallskip
{\parbox {6.5 cm} {Зав. кафедри \dotfill Ю.В.~Крак}}\hspace {0.7cm} {\hbox to 6.5 cm { Екзаменатор \dotfill Т.О.~Карнаух}} \end{scriptsize}}
%end
\vspace{0.9cm}
%begin
{{\begin {scriptsize}\par
{ \center  Київський національний університет імені Тараса Шевченка \par}
{\parbox {5 cm}{Освітня програма \hfill}{\parbox {7 cm}{Прикладна математика\hfill}}\parbox {6 cm}{\hfill Семестр 1}}\par
{\parbox {5 cm} {Навчальний предмет \hfill}\parbox {7 cm}{Програмування \hfill}\parbox {6cm}{\hfill}\par}\vspace{-15pt} \end{scriptsize}}{\center Екзаменаційний білет \No \textbf { 8 }\par}}
{%beginitem
1. Винятки та їх обробка.%enditem
\par
%beginitem
2. Написати програму, що запитує у користувача значення дійсних параметрів $x$ та $y$, обчислює для них значення виразу 
$$\frac{\log_{2} x + 18}{(\ctg x + 0.9) (y + 36)}$$ 
та повiдомляє введенi данi та результат обчислення.
 Оцінюється правильність та якість коду.

%enditem
\par
%beginitem
3. Написати функцію, що для цілого числа знаходить кількість його простих дільників, що входять у його розклад рівно у другому степеню. Наприклад, для числа $3^2{\cdot}7^9{\cdot}11{\cdot}23^{2}$ таких дільників два – $3$ та $23$, тому функція має повернути $2$.  Використання рядків та структур даних заборонено. Якість та швидкодія коду також оцінюються. 

%enditem
}
{\smallskip\smallskip \begin {scriptsize} Затверджено на засіданні кафедри теоретичної кібернетики, протокол \No 1  від   06.09.2023 р.\par\smallskip\smallskip\smallskip\smallskip
{\parbox {6.5 cm} {Зав. кафедри \dotfill Ю.В.~Крак}}\hspace {0.7cm} {\hbox to 6.5 cm { Екзаменатор \dotfill Т.О.~Карнаух}} \end{scriptsize}}
%end
\vspace{0.9cm}
%begin
{{\begin {scriptsize}\par
{ \center  Київський національний університет імені Тараса Шевченка \par}
{\parbox {5 cm}{Освітня програма \hfill}{\parbox {7 cm}{Прикладна математика\hfill}}\parbox {6 cm}{\hfill Семестр 1}}\par
{\parbox {5 cm} {Навчальний предмет \hfill}\parbox {7 cm}{Програмування \hfill}\parbox {6cm}{\hfill}\par}\vspace{-15pt} \end{scriptsize}}{\center Екзаменаційний білет \No \textbf { 9 }\par}}
{%beginitem
1. Рекурсивні означення, рекурсивні функції.%enditem
\par
%beginitem
2. Написати програму, що запитує у користувача значення дійсних параметрів $x$ та $y$, обчислює для них значення виразу 
$$\frac{\ln x + 52}{(\tg x + 0.2) (y + 49)}$$ 
та повiдомляє введенi данi та результат обчислення.
 Оцінюється правильність та якість коду.

%enditem
\par
%beginitem
3. Написати функцію, що для цілого числа знаходить простий дільник, який входить у розклад числа в найбільшому степені. Якщо таких кілька, то повернути найбільший. Наприклад, для числа $2^{9}{\cdot}3^5{\cdot}7^9{\cdot}11$ функція має повернути $7$.  Використання рядків та структур даних заборонено. Якість та швидкодія коду також оцінюються. 

%enditem
}
{\smallskip\smallskip \begin {scriptsize} Затверджено на засіданні кафедри теоретичної кібернетики, протокол \No 1  від   06.09.2023 р.\par\smallskip\smallskip\smallskip\smallskip
{\parbox {6.5 cm} {Зав. кафедри \dotfill Ю.В.~Крак}}\hspace {0.7cm} {\hbox to 6.5 cm { Екзаменатор \dotfill Т.О.~Карнаух}} \end{scriptsize}}
%end
\newpage
%begin
{{\begin {scriptsize}\par
{ \center  Київський національний університет імені Тараса Шевченка \par}
{\parbox {5 cm}{Освітня програма \hfill}{\parbox {7 cm}{Прикладна математика\hfill}}\parbox {6 cm}{\hfill Семестр 1}}\par
{\parbox {5 cm} {Навчальний предмет \hfill}\parbox {7 cm}{Програмування \hfill}\parbox {6cm}{\hfill}\par}\vspace{-15pt} \end{scriptsize}}{\center Екзаменаційний білет \No \textbf { 10 }\par}}
{%beginitem
1. Математичний підхід до визначення складності обчислень.%enditem
\par
%beginitem
2. Написати програму, що запитує у користувача значення дійсних параметрів $x$ та $y$, обчислює для них значення виразу 
$$\frac{\pi x - 19}{(\arcsin x - 0.3) (y - 39)}$$ 
та повiдомляє введенi данi та результат обчислення.
 Оцінюється правильність та якість коду.

%enditem
\par
%beginitem
3. Написати функцію, що для інтервалу $[a, b)$ знаходить знаходить число з найбільшою сумою простих дільників. Якщо таких чисел кілька, повертає найменше.  Використання рядків та структур даних заборонено. Якість та швидкодія коду також оцінюються. 

%enditem
}
{\smallskip\smallskip \begin {scriptsize} Затверджено на засіданні кафедри теоретичної кібернетики, протокол \No 1  від   06.09.2023 р.\par\smallskip\smallskip\smallskip\smallskip
{\parbox {6.5 cm} {Зав. кафедри \dotfill Ю.В.~Крак}}\hspace {0.7cm} {\hbox to 6.5 cm { Екзаменатор \dotfill Т.О.~Карнаух}} \end{scriptsize}}
%end
\vspace{0.9cm}
%begin
{{\begin {scriptsize}\par
{ \center  Київський національний університет імені Тараса Шевченка \par}
{\parbox {5 cm}{Освітня програма \hfill}{\parbox {7 cm}{Прикладна математика\hfill}}\parbox {6 cm}{\hfill Семестр 1}}\par
{\parbox {5 cm} {Навчальний предмет \hfill}\parbox {7 cm}{Програмування \hfill}\parbox {6cm}{\hfill}\par}\vspace{-15pt} \end{scriptsize}}{\center Екзаменаційний білет \No \textbf { 11 }\par}}
{%beginitem
1. Цикл while, умова продовження циклу, вигляд і семантика інструкції циклу while.%enditem
\par
%beginitem
2. Написати програму, що запитує у користувача значення дійсних параметрів $x$ та $y$, обчислює для них значення виразу 
$$\frac{\log_{10} x - 42}{(\arctg x + 0.5) (y + 80)}$$ 
та повiдомляє введенi данi та результат обчислення.
 Оцінюється правильність та якість коду.

%enditem
\par
%beginitem
3. Написати функцію, що для інтервалу $[a, b)$ знаходить кількість чисел, квадрати яких діляться на суму квадратів їх цифр. (Розглядаємо десятковий запис.)  Використання рядків та структур даних заборонено. Якість та швидкодія коду також оцінюються. 

%enditem
}
{\smallskip\smallskip \begin {scriptsize} Затверджено на засіданні кафедри теоретичної кібернетики, протокол \No 1  від   06.09.2023 р.\par\smallskip\smallskip\smallskip\smallskip
{\parbox {6.5 cm} {Зав. кафедри \dotfill Ю.В.~Крак}}\hspace {0.7cm} {\hbox to 6.5 cm { Екзаменатор \dotfill Т.О.~Карнаух}} \end{scriptsize}}
%end
\vspace{0.9cm}
%begin
{{\begin {scriptsize}\par
{ \center  Київський національний університет імені Тараса Шевченка \par}
{\parbox {5 cm}{Освітня програма \hfill}{\parbox {7 cm}{Прикладна математика\hfill}}\parbox {6 cm}{\hfill Семестр 1}}\par
{\parbox {5 cm} {Навчальний предмет \hfill}\parbox {7 cm}{Програмування \hfill}\parbox {6cm}{\hfill}\par}\vspace{-15pt} \end{scriptsize}}{\center Екзаменаційний білет \No \textbf { 12 }\par}}
{%beginitem
1. Глибина рекурсії та загальна кількість рекурсивних викликів, їх вплив на розміри програмного стеку та на час виконання програми.%enditem
\par
%beginitem
2. Написати програму, що запитує у користувача значення дійсних параметрів $x$ та $y$, обчислює для них значення виразу 
$$\frac{\ x + 59}{(\sin x - 0.1) (y - 91)}$$ 
та повiдомляє введенi данi та результат обчислення.
 Оцінюється правильність та якість коду.

%enditem
\par
%beginitem
3. Написати функцію, що для інтервалу $[a, b)$ знаходить найближчу пару чисел з цього інтервалу, що мають не менше трьох простих дільників та мають цифру 3 у своєму десятковому запису.  Використання рядків та структур даних заборонено. Якість та швидкодія коду також оцінюються. 

%enditem
}
{\smallskip\smallskip \begin {scriptsize} Затверджено на засіданні кафедри теоретичної кібернетики, протокол \No 1  від   06.09.2023 р.\par\smallskip\smallskip\smallskip\smallskip
{\parbox {6.5 cm} {Зав. кафедри \dotfill Ю.В.~Крак}}\hspace {0.7cm} {\hbox to 6.5 cm { Екзаменатор \dotfill Т.О.~Карнаух}} \end{scriptsize}}
%end
\newpage
%begin
{{\begin {scriptsize}\par
{ \center  Київський національний університет імені Тараса Шевченка \par}
{\parbox {5 cm}{Освітня програма \hfill}{\parbox {7 cm}{Прикладна математика\hfill}}\parbox {6 cm}{\hfill Семестр 1}}\par
{\parbox {5 cm} {Навчальний предмет \hfill}\parbox {7 cm}{Програмування \hfill}\parbox {6cm}{\hfill}\par}\vspace{-15pt} \end{scriptsize}}{\center Екзаменаційний білет \No \textbf { 13 }\par}}
{%beginitem
1. Емпіричний підхід до визначення складності обчислень.%enditem
\par
%beginitem
2. Написати програму, що запитує у користувача значення дійсних параметрів $x$ та $y$, обчислює для них значення виразу 
$$\frac{\log_{2} x - 74}{(\arccos x + 0.4) (y + 83)}$$ 
та повiдомляє введенi данi та результат обчислення.
 Оцінюється правильність та якість коду.

%enditem
\par
%beginitem
3. Написати функцію, що в інтервалі $[a, b)$ знаходить ціле число, що ділиться на найбільший степінь числа 3. Якщо таких чисел декілька, повернути найбільше.  Використання рядків та структур даних заборонено. Якість та швидкодія коду також оцінюються. 

%enditem
}
{\smallskip\smallskip \begin {scriptsize} Затверджено на засіданні кафедри теоретичної кібернетики, протокол \No 1  від   06.09.2023 р.\par\smallskip\smallskip\smallskip\smallskip
{\parbox {6.5 cm} {Зав. кафедри \dotfill Ю.В.~Крак}}\hspace {0.7cm} {\hbox to 6.5 cm { Екзаменатор \dotfill Т.О.~Карнаух}} \end{scriptsize}}
%end
\vspace{0.9cm}
%begin
{{\begin {scriptsize}\par
{ \center  Київський національний університет імені Тараса Шевченка \par}
{\parbox {5 cm}{Освітня програма \hfill}{\parbox {7 cm}{Прикладна математика\hfill}}\parbox {6 cm}{\hfill Семестр 1}}\par
{\parbox {5 cm} {Навчальний предмет \hfill}\parbox {7 cm}{Програмування \hfill}\parbox {6cm}{\hfill}\par}\vspace{-15pt} \end{scriptsize}}{\center Екзаменаційний білет \No \textbf { 14 }\par}}
{%beginitem
1. Арифметичні операції. Вирази. Пріоритети операторів. Порядок обчислення виразів.%enditem
\par
%beginitem
2. Написати програму, що запитує у користувача значення дійсних параметрів $x$ та $y$, обчислює для них значення виразу 
$$\frac{\ x + 95}{(\cos x - 0.7) (y - 25)}$$ 
та повiдомляє введенi данi та результат обчислення.
 Оцінюється правильність та якість коду.

%enditem
\par
%beginitem
3. Написати функцію, що для інтервалу $[a, b)$ знаходить найближчу пару чисел з цього інтервалу, що мають непарну кількість різних простих дільників. Якщо таких пар кілька, то цікавить пара з найменшим добутком.  Використання рядків та структур даних заборонено. Якість та швидкодія коду також оцінюються. 

%enditem
}
{\smallskip\smallskip \begin {scriptsize} Затверджено на засіданні кафедри теоретичної кібернетики, протокол \No 1  від   06.09.2023 р.\par\smallskip\smallskip\smallskip\smallskip
{\parbox {6.5 cm} {Зав. кафедри \dotfill Ю.В.~Крак}}\hspace {0.7cm} {\hbox to 6.5 cm { Екзаменатор \dotfill Т.О.~Карнаух}} \end{scriptsize}}
%end
\vspace{0.9cm}
%begin
{{\begin {scriptsize}\par
{ \center  Київський національний університет імені Тараса Шевченка \par}
{\parbox {5 cm}{Освітня програма \hfill}{\parbox {7 cm}{Прикладна математика\hfill}}\parbox {6 cm}{\hfill Семестр 1}}\par
{\parbox {5 cm} {Навчальний предмет \hfill}\parbox {7 cm}{Програмування \hfill}\parbox {6cm}{\hfill}\par}\vspace{-15pt} \end{scriptsize}}{\center Екзаменаційний білет \No \textbf { 15 }\par}}
{%beginitem
1. Змінні в Python. Типізація. Операції над змінними. Присвоювання та його види у Python.%enditem
\par
%beginitem
2. Написати програму, що запитує у користувача значення дійсних параметрів $x$ та $y$, обчислює для них значення виразу 
$$\frac{\pi x - 61}{(\arcsin x - 0.6) (y - 14)}$$ 
та повiдомляє введенi данi та результат обчислення.
 Оцінюється правильність та якість коду.

%enditem
\par
%beginitem
3. Написати функцію, що для інтервалу $[a, b)$ знаходить найближчу пару чисел з цього інтервалу, що мають парну кількість різних простих дільників. Якщо таких пар кілька, то цікавить пара з найменшим добутком.  Використання рядків та структур даних заборонено. Якість та швидкодія коду також оцінюються. 

%enditem
}
{\smallskip\smallskip \begin {scriptsize} Затверджено на засіданні кафедри теоретичної кібернетики, протокол \No 1  від   06.09.2023 р.\par\smallskip\smallskip\smallskip\smallskip
{\parbox {6.5 cm} {Зав. кафедри \dotfill Ю.В.~Крак}}\hspace {0.7cm} {\hbox to 6.5 cm { Екзаменатор \dotfill Т.О.~Карнаух}} \end{scriptsize}}
%end
\newpage
%begin
{{\begin {scriptsize}\par
{ \center  Київський національний університет імені Тараса Шевченка \par}
{\parbox {5 cm}{Освітня програма \hfill}{\parbox {7 cm}{Прикладна математика\hfill}}\parbox {6 cm}{\hfill Семестр 1}}\par
{\parbox {5 cm} {Навчальний предмет \hfill}\parbox {7 cm}{Програмування \hfill}\parbox {6cm}{\hfill}\par}\vspace{-15pt} \end{scriptsize}}{\center Екзаменаційний білет \No \textbf { 16 }\par}}
{%beginitem
1. Засоби керування порядком обчислень. Інструкція if.%enditem
\par
%beginitem
2. Написати програму, що запитує у користувача значення дійсних параметрів $x$ та $y$, обчислює для них значення виразу 
$$\frac{\log_{10} x + 38}{(\arctg x + 0.9) (y + 66)}$$ 
та повiдомляє введенi данi та результат обчислення.
 Оцінюється правильність та якість коду.

%enditem
\par
%beginitem
3. Написати функцію, яка для цілого числа знаходить суму його простих дільників, у десятковому запису яких нема цифри 3.  Використання рядків та структур даних заборонено. Якість та швидкодія коду також оцінюються. 

%enditem
}
{\smallskip\smallskip \begin {scriptsize} Затверджено на засіданні кафедри теоретичної кібернетики, протокол \No 1  від   06.09.2023 р.\par\smallskip\smallskip\smallskip\smallskip
{\parbox {6.5 cm} {Зав. кафедри \dotfill Ю.В.~Крак}}\hspace {0.7cm} {\hbox to 6.5 cm { Екзаменатор \dotfill Т.О.~Карнаух}} \end{scriptsize}}
%end
\vspace{0.9cm}
%begin
{{\begin {scriptsize}\par
{ \center  Київський національний університет імені Тараса Шевченка \par}
{\parbox {5 cm}{Освітня програма \hfill}{\parbox {7 cm}{Прикладна математика\hfill}}\parbox {6 cm}{\hfill Семестр 1}}\par
{\parbox {5 cm} {Навчальний предмет \hfill}\parbox {7 cm}{Програмування \hfill}\parbox {6cm}{\hfill}\par}\vspace{-15pt} \end{scriptsize}}{\center Екзаменаційний білет \No \textbf { 17 }\par}}
{%beginitem
1. Глобальні, локальні та вільні змінні.%enditem
\par
%beginitem
2. Написати програму, що запитує у користувача значення дійсних параметрів $x$ та $y$, обчислює для них значення виразу 
$$\frac{\sqrt x - 51}{(\cos x + 0.8) (y - 88)}$$ 
та повiдомляє введенi данi та результат обчислення.
 Оцінюється правильність та якість коду.

%enditem
\par
%beginitem
3. Написати функцію, що для цілого числа знаходить кількість простих дільників, які входять у розклад числа в найбільшому степені. Наприклад, для числа $3^5{\cdot}7^9{\cdot}11{\cdot}23^{9}$ таких дільників два – $23$ та $7$, тому функція має повернути $2$.  Використання рядків та структур даних заборонено. Якість та швидкодія коду також оцінюються. 

%enditem
}
{\smallskip\smallskip \begin {scriptsize} Затверджено на засіданні кафедри теоретичної кібернетики, протокол \No 1  від   06.09.2023 р.\par\smallskip\smallskip\smallskip\smallskip
{\parbox {6.5 cm} {Зав. кафедри \dotfill Ю.В.~Крак}}\hspace {0.7cm} {\hbox to 6.5 cm { Екзаменатор \dotfill Т.О.~Карнаух}} \end{scriptsize}}
%end
\vspace{0.9cm}
%begin
{{\begin {scriptsize}\par
{ \center  Київський національний університет імені Тараса Шевченка \par}
{\parbox {5 cm}{Освітня програма \hfill}{\parbox {7 cm}{Прикладна математика\hfill}}\parbox {6 cm}{\hfill Семестр 1}}\par
{\parbox {5 cm} {Навчальний предмет \hfill}\parbox {7 cm}{Програмування \hfill}\parbox {6cm}{\hfill}\par}\vspace{-15pt} \end{scriptsize}}{\center Екзаменаційний білет \No \textbf { 18 }\par}}
{%beginitem
1. Об’єктна система Python. Поняття незмінюваного об'єкту.%enditem
\par
%beginitem
2. Написати програму, що запитує у користувача значення дійсних параметрів $x$ та $y$, обчислює для них значення виразу 
$$\frac{\ln x + 40}{(\sin x - 0.2) (y + 69)}$$ 
та повiдомляє введенi данi та результат обчислення.
 Оцінюється правильність та якість коду.

%enditem
\par
%beginitem
3. Написати функцію, що для інтервалу $[a, b)$ знаходить найближчу пару чисел з цього інтервалу, що мають непарну кількість різних простих дільників. Якщо таких пар кілька, то цікавить пара з найбільшою сумою.  Використання рядків та структур даних заборонено. Якість та швидкодія коду також оцінюються. 

%enditem
}
{\smallskip\smallskip \begin {scriptsize} Затверджено на засіданні кафедри теоретичної кібернетики, протокол \No 1  від   06.09.2023 р.\par\smallskip\smallskip\smallskip\smallskip
{\parbox {6.5 cm} {Зав. кафедри \dotfill Ю.В.~Крак}}\hspace {0.7cm} {\hbox to 6.5 cm { Екзаменатор \dotfill Т.О.~Карнаух}} \end{scriptsize}}
%end
\newpage
%begin
{{\begin {scriptsize}\par
{ \center  Київський національний університет імені Тараса Шевченка \par}
{\parbox {5 cm}{Освітня програма \hfill}{\parbox {7 cm}{Прикладна математика\hfill}}\parbox {6 cm}{\hfill Семестр 1}}\par
{\parbox {5 cm} {Навчальний предмет \hfill}\parbox {7 cm}{Програмування \hfill}\parbox {6cm}{\hfill}\par}\vspace{-15pt} \end{scriptsize}}{\center Екзаменаційний білет \No \textbf { 19 }\par}}
{%beginitem
1. Компільовані та інтерпретовані мови.%enditem
\par
%beginitem
2. Написати програму, що запитує у користувача значення дійсних параметрів $x$ та $y$, обчислює для них значення виразу 
$$\frac{\log_{2} x - 70}{(\arccos x - 0.5) (y - 96)}$$ 
та повiдомляє введенi данi та результат обчислення.
 Оцінюється правильність та якість коду.

%enditem
\par
%beginitem
3. Написати функцію, що для інтервалу $[a, b)$ знаходить знаходить число з найменшою сумою простих дільників. Якщо таких чисел кілька, повертає найбільше.  Використання рядків та структур даних заборонено. Якість та швидкодія коду також оцінюються. 

%enditem
}
{\smallskip\smallskip \begin {scriptsize} Затверджено на засіданні кафедри теоретичної кібернетики, протокол \No 1  від   06.09.2023 р.\par\smallskip\smallskip\smallskip\smallskip
{\parbox {6.5 cm} {Зав. кафедри \dotfill Ю.В.~Крак}}\hspace {0.7cm} {\hbox to 6.5 cm { Екзаменатор \dotfill Т.О.~Карнаух}} \end{scriptsize}}
%end
\vspace{0.9cm}
%begin
{{\begin {scriptsize}\par
{ \center  Київський національний університет імені Тараса Шевченка \par}
{\parbox {5 cm}{Освітня програма \hfill}{\parbox {7 cm}{Прикладна математика\hfill}}\parbox {6 cm}{\hfill Семестр 1}}\par
{\parbox {5 cm} {Навчальний предмет \hfill}\parbox {7 cm}{Програмування \hfill}\parbox {6cm}{\hfill}\par}\vspace{-15pt} \end{scriptsize}}{\center Екзаменаційний білет \No \textbf { 20 }\par}}
{%beginitem
1. Компільовані та інтерпретовані мови.%enditem
\par
%beginitem
2. Написати програму, що запитує у користувача значення дійсних параметрів $x$ та $y$, обчислює для них значення виразу 
$$\frac{\ln x + 50}{(\tg x + 0.8) (y + 43)}$$ 
та повiдомляє введенi данi та результат обчислення.
 Оцінюється правильність та якість коду.

%enditem
\par
%beginitem
3. Написати функцію, що для інтервалу $[a, b)$ знаходить число з найбільшою кількістю різних простих дільників, що містять цифру 7 у своєму десятковому запису. Якщо їх кілька, то цікавить найменше.  Використання рядків та структур даних заборонено. Якість та швидкодія коду також оцінюються. 

%enditem
}
{\smallskip\smallskip \begin {scriptsize} Затверджено на засіданні кафедри теоретичної кібернетики, протокол \No 1  від   06.09.2023 р.\par\smallskip\smallskip\smallskip\smallskip
{\parbox {6.5 cm} {Зав. кафедри \dotfill Ю.В.~Крак}}\hspace {0.7cm} {\hbox to 6.5 cm { Екзаменатор \dotfill Т.О.~Карнаух}} \end{scriptsize}}
%end
\vspace{0.9cm}
%begin
{{\begin {scriptsize}\par
{ \center  Київський національний університет імені Тараса Шевченка \par}
{\parbox {5 cm}{Освітня програма \hfill}{\parbox {7 cm}{Прикладна математика\hfill}}\parbox {6 cm}{\hfill Семестр 1}}\par
{\parbox {5 cm} {Навчальний предмет \hfill}\parbox {7 cm}{Програмування \hfill}\parbox {6cm}{\hfill}\par}\vspace{-15pt} \end{scriptsize}}{\center Екзаменаційний білет \No \textbf { 21 }\par}}
{%beginitem
1. Лексеми мови Python.%enditem
\par
%beginitem
2. Написати програму, що запитує у користувача значення дійсних параметрів $x$ та $y$, обчислює для них значення виразу 
$$\frac{\log_{10} x + 41}{(\ctg x - 0.4) (y - 85)}$$ 
та повiдомляє введенi данi та результат обчислення.
 Оцінюється правильність та якість коду.

%enditem
\par
%beginitem
3. Написати функцію, що для цілого числа знаходить простий дільник, який входить у розклад числа в найбільшому степені. Якщо таких кілька, то повернути найменший. Наприклад, для числа $2^{9}{\cdot}3^5{\cdot}7^9{\cdot}11$ функція має повернути $2$.  Використання рядків та структур даних заборонено. Якість та швидкодія коду також оцінюються. 

%enditem
}
{\smallskip\smallskip \begin {scriptsize} Затверджено на засіданні кафедри теоретичної кібернетики, протокол \No 1  від   06.09.2023 р.\par\smallskip\smallskip\smallskip\smallskip
{\parbox {6.5 cm} {Зав. кафедри \dotfill Ю.В.~Крак}}\hspace {0.7cm} {\hbox to 6.5 cm { Екзаменатор \dotfill Т.О.~Карнаух}} \end{scriptsize}}
%end
\newpage
%begin
{{\begin {scriptsize}\par
{ \center  Київський національний університет імені Тараса Шевченка \par}
{\parbox {5 cm}{Освітня програма \hfill}{\parbox {7 cm}{Прикладна математика\hfill}}\parbox {6 cm}{\hfill Семестр 1}}\par
{\parbox {5 cm} {Навчальний предмет \hfill}\parbox {7 cm}{Програмування \hfill}\parbox {6cm}{\hfill}\par}\vspace{-15pt} \end{scriptsize}}{\center Екзаменаційний білет \No \textbf { 22 }\par}}
{%beginitem
1. Глобальні, локальні та вільні змінні.%enditem
\par
%beginitem
2. Написати програму, що запитує у користувача значення дійсних параметрів $x$ та $y$, обчислює для них значення виразу 
$$\frac{\ x - 47}{(\ctg x + 0.1) (y + 79)}$$ 
та повiдомляє введенi данi та результат обчислення.
 Оцінюється правильність та якість коду.

%enditem
\par
%beginitem
3. Написати функцію, що для інтервалу $[a, b)$ знаходить найближчу пару чисел з цього інтервалу, що мають від трьох до п'яти (включно) простих дільників.  Використання рядків та структур даних заборонено. Якість та швидкодія коду також оцінюються. 

%enditem
}
{\smallskip\smallskip \begin {scriptsize} Затверджено на засіданні кафедри теоретичної кібернетики, протокол \No 1  від   06.09.2023 р.\par\smallskip\smallskip\smallskip\smallskip
{\parbox {6.5 cm} {Зав. кафедри \dotfill Ю.В.~Крак}}\hspace {0.7cm} {\hbox to 6.5 cm { Екзаменатор \dotfill Т.О.~Карнаух}} \end{scriptsize}}
%end
\vspace{0.9cm}
%begin
{{\begin {scriptsize}\par
{ \center  Київський національний університет імені Тараса Шевченка \par}
{\parbox {5 cm}{Освітня програма \hfill}{\parbox {7 cm}{Прикладна математика\hfill}}\parbox {6 cm}{\hfill Семестр 1}}\par
{\parbox {5 cm} {Навчальний предмет \hfill}\parbox {7 cm}{Програмування \hfill}\parbox {6cm}{\hfill}\par}\vspace{-15pt} \end{scriptsize}}{\center Екзаменаційний білет \No \textbf { 23 }\par}}
{%beginitem
1. Цикл while, умова продовження циклу, вигляд і семантика інструкції циклу while.%enditem
\par
%beginitem
2. Написати програму, що запитує у користувача значення дійсних параметрів $x$ та $y$, обчислює для них значення виразу 
$$\frac{\pi x + 57}{(\arccos x + 0.9) (y - 98)}$$ 
та повiдомляє введенi данi та результат обчислення.
 Оцінюється правильність та якість коду.

%enditem
\par
%beginitem
3. Написати функцію, що для цілого числа знаходить кількість його простих дільників, що входять у його розклад не більш ніж у 5му степеню. Наприклад, для числа $3^8{\cdot}7^9{\cdot}11{\cdot}23^{2}$ таких дільників два – $11$ та $23$, тому функція має повернути $2$.  Використання рядків та структур даних заборонено. Якість та швидкодія коду також оцінюються. 

%enditem
}
{\smallskip\smallskip \begin {scriptsize} Затверджено на засіданні кафедри теоретичної кібернетики, протокол \No 1  від   06.09.2023 р.\par\smallskip\smallskip\smallskip\smallskip
{\parbox {6.5 cm} {Зав. кафедри \dotfill Ю.В.~Крак}}\hspace {0.7cm} {\hbox to 6.5 cm { Екзаменатор \dotfill Т.О.~Карнаух}} \end{scriptsize}}
%end
\vspace{0.9cm}
%begin
{{\begin {scriptsize}\par
{ \center  Київський національний університет імені Тараса Шевченка \par}
{\parbox {5 cm}{Освітня програма \hfill}{\parbox {7 cm}{Прикладна математика\hfill}}\parbox {6 cm}{\hfill Семестр 1}}\par
{\parbox {5 cm} {Навчальний предмет \hfill}\parbox {7 cm}{Програмування \hfill}\parbox {6cm}{\hfill}\par}\vspace{-15pt} \end{scriptsize}}{\center Екзаменаційний білет \No \textbf { 24 }\par}}
{%beginitem
1. Вбудовані числові типи мови Python.%enditem
\par
%beginitem
2. Написати програму, що запитує у користувача значення дійсних параметрів $x$ та $y$, обчислює для них значення виразу 
$$\frac{\sqrt x - 33}{(\tg x - 0.3) (y + 72)}$$ 
та повiдомляє введенi данi та результат обчислення.
 Оцінюється правильність та якість коду.

%enditem
\par
%beginitem
3. Написати функцію, що в інтервалі $[a, b)$ знаходить ціле число, що ділиться на найбільший степінь числа 3. Якщо таких чисел декілька, повернути найменше.  Використання рядків та структур даних заборонено. Якість та швидкодія коду також оцінюються. 

%enditem
}
{\smallskip\smallskip \begin {scriptsize} Затверджено на засіданні кафедри теоретичної кібернетики, протокол \No 1  від   06.09.2023 р.\par\smallskip\smallskip\smallskip\smallskip
{\parbox {6.5 cm} {Зав. кафедри \dotfill Ю.В.~Крак}}\hspace {0.7cm} {\hbox to 6.5 cm { Екзаменатор \dotfill Т.О.~Карнаух}} \end{scriptsize}}
%end
\newpage
%begin
{{\begin {scriptsize}\par
{ \center  Київський національний університет імені Тараса Шевченка \par}
{\parbox {5 cm}{Освітня програма \hfill}{\parbox {7 cm}{Прикладна математика\hfill}}\parbox {6 cm}{\hfill Семестр 1}}\par
{\parbox {5 cm} {Навчальний предмет \hfill}\parbox {7 cm}{Програмування \hfill}\parbox {6cm}{\hfill}\par}\vspace{-15pt} \end{scriptsize}}{\center Екзаменаційний білет \No \textbf { 25 }\par}}
{%beginitem
1. Арифметичні операції. Вирази. Пріоритети операторів. Порядок обчислення виразів.%enditem
\par
%beginitem
2. Написати програму, що запитує у користувача значення дійсних параметрів $x$ та $y$, обчислює для них значення виразу 
$$\frac{\pi x - 27}{(\sin x - 0.2) (y + 16)}$$ 
та повiдомляє введенi данi та результат обчислення.
 Оцінюється правильність та якість коду.

%enditem
\par
%beginitem
3. Написати функцію, що для інтервалу $[a, b)$ знаходить число з найбільшою кількістю різних простих дільників. Якщо їх кілька, то цікавить найбільше.  Використання рядків та структур даних заборонено. Якість та швидкодія коду також оцінюються. 

%enditem
}
{\smallskip\smallskip \begin {scriptsize} Затверджено на засіданні кафедри теоретичної кібернетики, протокол \No 1  від   06.09.2023 р.\par\smallskip\smallskip\smallskip\smallskip
{\parbox {6.5 cm} {Зав. кафедри \dotfill Ю.В.~Крак}}\hspace {0.7cm} {\hbox to 6.5 cm { Екзаменатор \dotfill Т.О.~Карнаух}} \end{scriptsize}}
%end
\vspace{0.9cm}
%begin
{{\begin {scriptsize}\par
{ \center  Київський національний університет імені Тараса Шевченка \par}
{\parbox {5 cm}{Освітня програма \hfill}{\parbox {7 cm}{Прикладна математика\hfill}}\parbox {6 cm}{\hfill Семестр 1}}\par
{\parbox {5 cm} {Навчальний предмет \hfill}\parbox {7 cm}{Програмування \hfill}\parbox {6cm}{\hfill}\par}\vspace{-15pt} \end{scriptsize}}{\center Екзаменаційний білет \No \textbf { 26 }\par}}
{%beginitem
1. Поняття часової та просторової складності обчислень. Складність у середньому. Асимптотичні оцінки складності.%enditem
\par
%beginitem
2. Написати програму, що запитує у користувача значення дійсних параметрів $x$ та $y$, обчислює для них значення виразу 
$$\frac{\sqrt x + 31}{(\arctg x + 0.6) (y - 63)}$$ 
та повiдомляє введенi данi та результат обчислення.
 Оцінюється правильність та якість коду.

%enditem
\par
%beginitem
3. Написати функцію, що для цілого числа знаходить суму його простих дільників, що входять у його розклад не більш ніж у 4му степеню. Наприклад, для числа $3^5{\cdot}7^9{\cdot}11{\cdot}23^{2}$ таких дільників два – $11$ та $23$, тому функція має повернути $34$.  Використання рядків та структур даних заборонено. Якість та швидкодія коду також оцінюються. 

%enditem
}
{\smallskip\smallskip \begin {scriptsize} Затверджено на засіданні кафедри теоретичної кібернетики, протокол \No 1  від   06.09.2023 р.\par\smallskip\smallskip\smallskip\smallskip
{\parbox {6.5 cm} {Зав. кафедри \dotfill Ю.В.~Крак}}\hspace {0.7cm} {\hbox to 6.5 cm { Екзаменатор \dotfill Т.О.~Карнаух}} \end{scriptsize}}
%end
\vspace{0.9cm}
%begin
{{\begin {scriptsize}\par
{ \center  Київський національний університет імені Тараса Шевченка \par}
{\parbox {5 cm}{Освітня програма \hfill}{\parbox {7 cm}{Прикладна математика\hfill}}\parbox {6 cm}{\hfill Семестр 1}}\par
{\parbox {5 cm} {Навчальний предмет \hfill}\parbox {7 cm}{Програмування \hfill}\parbox {6cm}{\hfill}\par}\vspace{-15pt} \end{scriptsize}}{\center Екзаменаційний білет \No \textbf { 27 }\par}}
{%beginitem
1. Емпіричний підхід до визначення складності обчислень.%enditem
\par
%beginitem
2. Написати програму, що запитує у користувача значення дійсних параметрів $x$ та $y$, обчислює для них значення виразу 
$$\frac{\log_{10} x - 90}{(\cos x - 0.7) (y + 87)}$$ 
та повiдомляє введенi данi та результат обчислення.
 Оцінюється правильність та якість коду.

%enditem
\par
%beginitem
3. Написати функцію, що для інтервалу $[a, b)$ знаходить знаходить число з найбільшою сумою простих дільників. Якщо таких чисел кілька, повертає найменше.  Використання рядків та структур даних заборонено. Якість та швидкодія коду також оцінюються. 

%enditem
}
{\smallskip\smallskip \begin {scriptsize} Затверджено на засіданні кафедри теоретичної кібернетики, протокол \No 1  від   06.09.2023 р.\par\smallskip\smallskip\smallskip\smallskip
{\parbox {6.5 cm} {Зав. кафедри \dotfill Ю.В.~Крак}}\hspace {0.7cm} {\hbox to 6.5 cm { Екзаменатор \dotfill Т.О.~Карнаух}} \end{scriptsize}}
%end
\newpage
%begin
{{\begin {scriptsize}\par
{ \center  Київський національний університет імені Тараса Шевченка \par}
{\parbox {5 cm}{Освітня програма \hfill}{\parbox {7 cm}{Прикладна математика\hfill}}\parbox {6 cm}{\hfill Семестр 1}}\par
{\parbox {5 cm} {Навчальний предмет \hfill}\parbox {7 cm}{Програмування \hfill}\parbox {6cm}{\hfill}\par}\vspace{-15pt} \end{scriptsize}}{\center Екзаменаційний білет \No \textbf { 28 }\par}}
{%beginitem
1. Рекурсивні означення, рекурсивні функції.%enditem
\par
%beginitem
2. Написати програму, що запитує у користувача значення дійсних параметрів $x$ та $y$, обчислює для них значення виразу 
$$\frac{\ x + 22}{(\arcsin x + 0.8) (y - 76)}$$ 
та повiдомляє введенi данi та результат обчислення.
 Оцінюється правильність та якість коду.

%enditem
\par
%beginitem
3. Написати функцію, що для інтервалу $[a, b)$ знаходить знаходить число з найменшою сумою простих дільників. Якщо таких чисел кілька, повертає найбільше.  Використання рядків та структур даних заборонено. Якість та швидкодія коду також оцінюються. 

%enditem
}
{\smallskip\smallskip \begin {scriptsize} Затверджено на засіданні кафедри теоретичної кібернетики, протокол \No 1  від   06.09.2023 р.\par\smallskip\smallskip\smallskip\smallskip
{\parbox {6.5 cm} {Зав. кафедри \dotfill Ю.В.~Крак}}\hspace {0.7cm} {\hbox to 6.5 cm { Екзаменатор \dotfill Т.О.~Карнаух}} \end{scriptsize}}
%end
\vspace{0.9cm}
%begin
{{\begin {scriptsize}\par
{ \center  Київський національний університет імені Тараса Шевченка \par}
{\parbox {5 cm}{Освітня програма \hfill}{\parbox {7 cm}{Прикладна математика\hfill}}\parbox {6 cm}{\hfill Семестр 1}}\par
{\parbox {5 cm} {Навчальний предмет \hfill}\parbox {7 cm}{Програмування \hfill}\parbox {6cm}{\hfill}\par}\vspace{-15pt} \end{scriptsize}}{\center Екзаменаційний білет \No \textbf { 29 }\par}}
{%beginitem
1. Цикл for, вигляд і семантика інструкції циклу for.%enditem
\par
%beginitem
2. Написати програму, що запитує у користувача значення дійсних параметрів $x$ та $y$, обчислює для них значення виразу 
$$\frac{\ln x + 13}{(\cos x - 0.1) (y - 35)}$$ 
та повiдомляє введенi данi та результат обчислення.
 Оцінюється правильність та якість коду.

%enditem
\par
%beginitem
3. Написати функцію, що для інтервалу $[a, b)$ знаходить найближчу пару чисел з цього інтервалу, що мають від трьох до п'яти (включно) простих дільників.  Використання рядків та структур даних заборонено. Якість та швидкодія коду також оцінюються. 

%enditem
}
{\smallskip\smallskip \begin {scriptsize} Затверджено на засіданні кафедри теоретичної кібернетики, протокол \No 1  від   06.09.2023 р.\par\smallskip\smallskip\smallskip\smallskip
{\parbox {6.5 cm} {Зав. кафедри \dotfill Ю.В.~Крак}}\hspace {0.7cm} {\hbox to 6.5 cm { Екзаменатор \dotfill Т.О.~Карнаух}} \end{scriptsize}}
%end
\vspace{0.9cm}
%begin
{{\begin {scriptsize}\par
{ \center  Київський національний університет імені Тараса Шевченка \par}
{\parbox {5 cm}{Освітня програма \hfill}{\parbox {7 cm}{Прикладна математика\hfill}}\parbox {6 cm}{\hfill Семестр 1}}\par
{\parbox {5 cm} {Навчальний предмет \hfill}\parbox {7 cm}{Програмування \hfill}\parbox {6cm}{\hfill}\par}\vspace{-15pt} \end{scriptsize}}{\center Екзаменаційний білет \No \textbf { 30 }\par}}
{%beginitem
1. Змінні в Python. Типізація. Операції над змінними. Присвоювання та його види у Python.%enditem
\par
%beginitem
2. Написати програму, що запитує у користувача значення дійсних параметрів $x$ та $y$, обчислює для них значення виразу 
$$\frac{\log_{2} x - 93}{(\ctg x + 0.7) (y + 97)}$$ 
та повiдомляє введенi данi та результат обчислення.
 Оцінюється правильність та якість коду.

%enditem
\par
%beginitem
3. Написати функцію, що для інтервалу $[a, b)$ знаходить знаходить число з найменшим добутком простих дільників. Якщо таких чисел кілька, повертає найбільше.  Використання рядків та структур даних заборонено. Якість та швидкодія коду також оцінюються. 

%enditem
}
{\smallskip\smallskip \begin {scriptsize} Затверджено на засіданні кафедри теоретичної кібернетики, протокол \No 1  від   06.09.2023 р.\par\smallskip\smallskip\smallskip\smallskip
{\parbox {6.5 cm} {Зав. кафедри \dotfill Ю.В.~Крак}}\hspace {0.7cm} {\hbox to 6.5 cm { Екзаменатор \dotfill Т.О.~Карнаух}} \end{scriptsize}}
%end
\newpage
%begin
{{\begin {scriptsize}\par
{ \center  Київський національний університет імені Тараса Шевченка \par}
{\parbox {5 cm}{Освітня програма \hfill}{\parbox {7 cm}{Прикладна математика\hfill}}\parbox {6 cm}{\hfill Семестр 1}}\par
{\parbox {5 cm} {Навчальний предмет \hfill}\parbox {7 cm}{Програмування \hfill}\parbox {6cm}{\hfill}\par}\vspace{-15pt} \end{scriptsize}}{\center Екзаменаційний білет \No \textbf { 31 }\par}}
{%beginitem
1. Глибина рекурсії та загальна кількість рекурсивних викликів, їх вплив на розміри програмного стеку та на час виконання програми.%enditem
\par
%beginitem
2. Написати програму, що запитує у користувача значення дійсних параметрів $x$ та $y$, обчислює для них значення виразу 
$$\frac{\ x - 46}{(\tg x + 0.3) (y + 55)}$$ 
та повiдомляє введенi данi та результат обчислення.
 Оцінюється правильність та якість коду.

%enditem
\par
%beginitem
3. Написати функцію, що для інтервалу $[a, b)$ знаходить знаходить найбільший степінь числа 3, на який ділиться хоча б одне ціле число з цього інтервалу.  Використання рядків та структур даних заборонено. Якість та швидкодія коду також оцінюються. 

%enditem
}
{\smallskip\smallskip \begin {scriptsize} Затверджено на засіданні кафедри теоретичної кібернетики, протокол \No 1  від   06.09.2023 р.\par\smallskip\smallskip\smallskip\smallskip
{\parbox {6.5 cm} {Зав. кафедри \dotfill Ю.В.~Крак}}\hspace {0.7cm} {\hbox to 6.5 cm { Екзаменатор \dotfill Т.О.~Карнаух}} \end{scriptsize}}
%end
\vspace{0.9cm}
%begin
{{\begin {scriptsize}\par
{ \center  Київський національний університет імені Тараса Шевченка \par}
{\parbox {5 cm}{Освітня програма \hfill}{\parbox {7 cm}{Прикладна математика\hfill}}\parbox {6 cm}{\hfill Семестр 1}}\par
{\parbox {5 cm} {Навчальний предмет \hfill}\parbox {7 cm}{Програмування \hfill}\parbox {6cm}{\hfill}\par}\vspace{-15pt} \end{scriptsize}}{\center Екзаменаційний білет \No \textbf { 32 }\par}}
{%beginitem
1. Оператори порівняння (у широкому розумінні мови Python) та булеві оператори.%enditem
\par
%beginitem
2. Написати програму, що запитує у користувача значення дійсних параметрів $x$ та $y$, обчислює для них значення виразу 
$$\frac{\log_{2} x + 29}{(\arccos x - 0.6) (y - 28)}$$ 
та повiдомляє введенi данi та результат обчислення.
 Оцінюється правильність та якість коду.

%enditem
\par
%beginitem
3. Написати функцію, яка для цілого числа знаходить кількість його різних простих дільників, що мають у десятковому запису цифру 1.  Використання рядків та структур даних заборонено. Якість та швидкодія коду також оцінюються. 

%enditem
}
{\smallskip\smallskip \begin {scriptsize} Затверджено на засіданні кафедри теоретичної кібернетики, протокол \No 1  від   06.09.2023 р.\par\smallskip\smallskip\smallskip\smallskip
{\parbox {6.5 cm} {Зав. кафедри \dotfill Ю.В.~Крак}}\hspace {0.7cm} {\hbox to 6.5 cm { Екзаменатор \dotfill Т.О.~Карнаух}} \end{scriptsize}}
%end
\vspace{0.9cm}
%begin
{{\begin {scriptsize}\par
{ \center  Київський національний університет імені Тараса Шевченка \par}
{\parbox {5 cm}{Освітня програма \hfill}{\parbox {7 cm}{Прикладна математика\hfill}}\parbox {6 cm}{\hfill Семестр 1}}\par
{\parbox {5 cm} {Навчальний предмет \hfill}\parbox {7 cm}{Програмування \hfill}\parbox {6cm}{\hfill}\par}\vspace{-15pt} \end{scriptsize}}{\center Екзаменаційний білет \No \textbf { 33 }\par}}
{%beginitem
1. Винятки та їх обробка.%enditem
\par
%beginitem
2. Написати програму, що запитує у користувача значення дійсних параметрів $x$ та $y$, обчислює для них значення виразу 
$$\frac{\log_{10} x + 53}{(\arctg x + 0.4) (y + 99)}$$ 
та повiдомляє введенi данi та результат обчислення.
 Оцінюється правильність та якість коду.

%enditem
\par
%beginitem
3. Написати функцію, що для інтервалу $[a, b)$ знаходить найближчу пару чисел з цього інтервалу, що мають непарну кількість різних простих дільників. Якщо таких пар кілька, то цікавить пара з найбільшою сумою.  Використання рядків та структур даних заборонено. Якість та швидкодія коду також оцінюються. 

%enditem
}
{\smallskip\smallskip \begin {scriptsize} Затверджено на засіданні кафедри теоретичної кібернетики, протокол \No 1  від   06.09.2023 р.\par\smallskip\smallskip\smallskip\smallskip
{\parbox {6.5 cm} {Зав. кафедри \dotfill Ю.В.~Крак}}\hspace {0.7cm} {\hbox to 6.5 cm { Екзаменатор \dotfill Т.О.~Карнаух}} \end{scriptsize}}
%end
\newpage
%begin
{{\begin {scriptsize}\par
{ \center  Київський національний університет імені Тараса Шевченка \par}
{\parbox {5 cm}{Освітня програма \hfill}{\parbox {7 cm}{Прикладна математика\hfill}}\parbox {6 cm}{\hfill Семестр 1}}\par
{\parbox {5 cm} {Навчальний предмет \hfill}\parbox {7 cm}{Програмування \hfill}\parbox {6cm}{\hfill}\par}\vspace{-15pt} \end{scriptsize}}{\center Екзаменаційний білет \No \textbf { 34 }\par}}
{%beginitem
1. Поняття типу даних.%enditem
\par
%beginitem
2. Написати програму, що запитує у користувача значення дійсних параметрів $x$ та $y$, обчислює для них значення виразу 
$$\frac{\ln x - 58}{(\arcsin x - 0.5) (y - 56)}$$ 
та повiдомляє введенi данi та результат обчислення.
 Оцінюється правильність та якість коду.

%enditem
\par
%beginitem
3. Написати функцію, що для інтервалу $[a, b)$ знаходить число з найбільшою кількістю різних простих дільників, що містять цифру 7 у своєму десятковому запису. Якщо їх кілька, то цікавить найбільше.  Використання рядків та структур даних заборонено. Якість та швидкодія коду також оцінюються. 

%enditem
}
{\smallskip\smallskip \begin {scriptsize} Затверджено на засіданні кафедри теоретичної кібернетики, протокол \No 1  від   06.09.2023 р.\par\smallskip\smallskip\smallskip\smallskip
{\parbox {6.5 cm} {Зав. кафедри \dotfill Ю.В.~Крак}}\hspace {0.7cm} {\hbox to 6.5 cm { Екзаменатор \dotfill Т.О.~Карнаух}} \end{scriptsize}}
%end
\vspace{0.9cm}
%begin
{{\begin {scriptsize}\par
{ \center  Київський національний університет імені Тараса Шевченка \par}
{\parbox {5 cm}{Освітня програма \hfill}{\parbox {7 cm}{Прикладна математика\hfill}}\parbox {6 cm}{\hfill Семестр 1}}\par
{\parbox {5 cm} {Навчальний предмет \hfill}\parbox {7 cm}{Програмування \hfill}\parbox {6cm}{\hfill}\par}\vspace{-15pt} \end{scriptsize}}{\center Екзаменаційний білет \No \textbf { 35 }\par}}
{%beginitem
1. Об’єктна система Python. Поняття незмінюваного об'єкту.%enditem
\par
%beginitem
2. Написати програму, що запитує у користувача значення дійсних параметрів $x$ та $y$, обчислює для них значення виразу 
$$\frac{\sqrt x + 26}{(\sin x - 0.9) (y + 65)}$$ 
та повiдомляє введенi данi та результат обчислення.
 Оцінюється правильність та якість коду.

%enditem
\par
%beginitem
3. Написати функцію, що в інтервалі $[a, b)$ знаходить ціле число, що ділиться на найбільший степінь числа 3. Якщо таких чисел декілька, повернути найбільше.  Використання рядків та структур даних заборонено. Якість та швидкодія коду також оцінюються. 

%enditem
}
{\smallskip\smallskip \begin {scriptsize} Затверджено на засіданні кафедри теоретичної кібернетики, протокол \No 1  від   06.09.2023 р.\par\smallskip\smallskip\smallskip\smallskip
{\parbox {6.5 cm} {Зав. кафедри \dotfill Ю.В.~Крак}}\hspace {0.7cm} {\hbox to 6.5 cm { Екзаменатор \dotfill Т.О.~Карнаух}} \end{scriptsize}}
%end
\vspace{0.9cm}
%begin
{{\begin {scriptsize}\par
{ \center  Київський національний університет імені Тараса Шевченка \par}
{\parbox {5 cm}{Освітня програма \hfill}{\parbox {7 cm}{Прикладна математика\hfill}}\parbox {6 cm}{\hfill Семестр 1}}\par
{\parbox {5 cm} {Навчальний предмет \hfill}\parbox {7 cm}{Програмування \hfill}\parbox {6cm}{\hfill}\par}\vspace{-15pt} \end{scriptsize}}{\center Екзаменаційний білет \No \textbf { 36 }\par}}
{%beginitem
1. Математичний підхід до визначення складності обчислень.%enditem
\par
%beginitem
2. Написати програму, що запитує у користувача значення дійсних параметрів $x$ та $y$, обчислює для них значення виразу 
$$\frac{\pi x - 20}{(\cos x + 0.2) (y - 15)}$$ 
та повiдомляє введенi данi та результат обчислення.
 Оцінюється правильність та якість коду.

%enditem
\par
%beginitem
3. Написати функцію, що для цілого числа знаходить суму його простих дільників, що входять у його розклад не більш ніж у 4му степеню. Наприклад, для числа $3^5{\cdot}7^9{\cdot}11{\cdot}23^{2}$ таких дільників два – $11$ та $23$, тому функція має повернути $34$.  Використання рядків та структур даних заборонено. Якість та швидкодія коду також оцінюються. 

%enditem
}
{\smallskip\smallskip \begin {scriptsize} Затверджено на засіданні кафедри теоретичної кібернетики, протокол \No 1  від   06.09.2023 р.\par\smallskip\smallskip\smallskip\smallskip
{\parbox {6.5 cm} {Зав. кафедри \dotfill Ю.В.~Крак}}\hspace {0.7cm} {\hbox to 6.5 cm { Екзаменатор \dotfill Т.О.~Карнаух}} \end{scriptsize}}
%end
\newpage
%begin
{{\begin {scriptsize}\par
{ \center  Київський національний університет імені Тараса Шевченка \par}
{\parbox {5 cm}{Освітня програма \hfill}{\parbox {7 cm}{Прикладна математика\hfill}}\parbox {6 cm}{\hfill Семестр 1}}\par
{\parbox {5 cm} {Навчальний предмет \hfill}\parbox {7 cm}{Програмування \hfill}\parbox {6cm}{\hfill}\par}\vspace{-15pt} \end{scriptsize}}{\center Екзаменаційний білет \No \textbf { 37 }\par}}
{%beginitem
1. Семантика функції та її виклику, параметри та аргументи.%enditem
\par
%beginitem
2. Написати програму, що запитує у користувача значення дійсних параметрів $x$ та $y$, обчислює для них значення виразу 
$$\frac{\ x - 24}{(\ctg x + 0.5) (y + 48)}$$ 
та повiдомляє введенi данi та результат обчислення.
 Оцінюється правильність та якість коду.

%enditem
\par
%beginitem
3. Написати функцію, що в інтервалі $[a, b)$ знаходить ціле число, що ділиться на найбільший степінь числа 3. Якщо таких чисел декілька, повернути найменше.  Використання рядків та структур даних заборонено. Якість та швидкодія коду також оцінюються. 

%enditem
}
{\smallskip\smallskip \begin {scriptsize} Затверджено на засіданні кафедри теоретичної кібернетики, протокол \No 1  від   06.09.2023 р.\par\smallskip\smallskip\smallskip\smallskip
{\parbox {6.5 cm} {Зав. кафедри \dotfill Ю.В.~Крак}}\hspace {0.7cm} {\hbox to 6.5 cm { Екзаменатор \dotfill Т.О.~Карнаух}} \end{scriptsize}}
%end
\vspace{0.9cm}
%begin
{{\begin {scriptsize}\par
{ \center  Київський національний університет імені Тараса Шевченка \par}
{\parbox {5 cm}{Освітня програма \hfill}{\parbox {7 cm}{Прикладна математика\hfill}}\parbox {6 cm}{\hfill Семестр 1}}\par
{\parbox {5 cm} {Навчальний предмет \hfill}\parbox {7 cm}{Програмування \hfill}\parbox {6cm}{\hfill}\par}\vspace{-15pt} \end{scriptsize}}{\center Екзаменаційний білет \No \textbf { 38 }\par}}
{%beginitem
1. Засоби керування порядком обчислень. Інструкція if.%enditem
\par
%beginitem
2. Написати програму, що запитує у користувача значення дійсних параметрів $x$ та $y$, обчислює для них значення виразу 
$$\frac{\sqrt x + 86}{(\arcsin x - 0.2) (y - 73)}$$ 
та повiдомляє введенi данi та результат обчислення.
 Оцінюється правильність та якість коду.

%enditem
\par
%beginitem
3. Написати функцію, що для інтервалу $[a, b)$ знаходить найближчу пару чисел з цього інтервалу, що мають не менше трьох простих дільників та мають цифру 3 у своєму десятковому запису.  Використання рядків та структур даних заборонено. Якість та швидкодія коду також оцінюються. 

%enditem
}
{\smallskip\smallskip \begin {scriptsize} Затверджено на засіданні кафедри теоретичної кібернетики, протокол \No 1  від   06.09.2023 р.\par\smallskip\smallskip\smallskip\smallskip
{\parbox {6.5 cm} {Зав. кафедри \dotfill Ю.В.~Крак}}\hspace {0.7cm} {\hbox to 6.5 cm { Екзаменатор \dotfill Т.О.~Карнаух}} \end{scriptsize}}
%end
\vspace{0.9cm}
%begin
{{\begin {scriptsize}\par
{ \center  Київський національний університет імені Тараса Шевченка \par}
{\parbox {5 cm}{Освітня програма \hfill}{\parbox {7 cm}{Прикладна математика\hfill}}\parbox {6 cm}{\hfill Семестр 1}}\par
{\parbox {5 cm} {Навчальний предмет \hfill}\parbox {7 cm}{Програмування \hfill}\parbox {6cm}{\hfill}\par}\vspace{-15pt} \end{scriptsize}}{\center Екзаменаційний білет \No \textbf { 39 }\par}}
{%beginitem
1. Цикл for, вигляд і семантика інструкції циклу for.%enditem
\par
%beginitem
2. Написати програму, що запитує у користувача значення дійсних параметрів $x$ та $y$, обчислює для них значення виразу 
$$\frac{\log_{2} x - 71}{(\sin x + 0.1) (y - 75)}$$ 
та повiдомляє введенi данi та результат обчислення.
 Оцінюється правильність та якість коду.

%enditem
\par
%beginitem
3. Написати функцію, що для інтервалу $[a, b)$ знаходить кількість чисел, квадрати яких діляться на суму квадратів їх цифр. (Розглядаємо десятковий запис.)  Використання рядків та структур даних заборонено. Якість та швидкодія коду також оцінюються. 

%enditem
}
{\smallskip\smallskip \begin {scriptsize} Затверджено на засіданні кафедри теоретичної кібернетики, протокол \No 1  від   06.09.2023 р.\par\smallskip\smallskip\smallskip\smallskip
{\parbox {6.5 cm} {Зав. кафедри \dotfill Ю.В.~Крак}}\hspace {0.7cm} {\hbox to 6.5 cm { Екзаменатор \dotfill Т.О.~Карнаух}} \end{scriptsize}}
%end
\newpage
%begin
{{\begin {scriptsize}\par
{ \center  Київський національний університет імені Тараса Шевченка \par}
{\parbox {5 cm}{Освітня програма \hfill}{\parbox {7 cm}{Прикладна математика\hfill}}\parbox {6 cm}{\hfill Семестр 1}}\par
{\parbox {5 cm} {Навчальний предмет \hfill}\parbox {7 cm}{Програмування \hfill}\parbox {6cm}{\hfill}\par}\vspace{-15pt} \end{scriptsize}}{\center Екзаменаційний білет \No \textbf { 40 }\par}}
{%beginitem
1. Семантика функції та її виклику, параметри та аргументи.%enditem
\par
%beginitem
2. Написати програму, що запитує у користувача значення дійсних параметрів $x$ та $y$, обчислює для них значення виразу 
$$\frac{\pi x + 94}{(\arctg x - 0.7) (y + 17)}$$ 
та повiдомляє введенi данi та результат обчислення.
 Оцінюється правильність та якість коду.

%enditem
\par
%beginitem
3. Написати функцію, що для цілого числа знаходить кількість його простих дільників, що входять у його розклад не більш ніж у 5му степеню. Наприклад, для числа $3^8{\cdot}7^9{\cdot}11{\cdot}23^{2}$ таких дільників два – $11$ та $23$, тому функція має повернути $2$.  Використання рядків та структур даних заборонено. Якість та швидкодія коду також оцінюються. 

%enditem
}
{\smallskip\smallskip \begin {scriptsize} Затверджено на засіданні кафедри теоретичної кібернетики, протокол \No 1  від   06.09.2023 р.\par\smallskip\smallskip\smallskip\smallskip
{\parbox {6.5 cm} {Зав. кафедри \dotfill Ю.В.~Крак}}\hspace {0.7cm} {\hbox to 6.5 cm { Екзаменатор \dotfill Т.О.~Карнаух}} \end{scriptsize}}
%end
\vspace{0.9cm}
%begin
{{\begin {scriptsize}\par
{ \center  Київський національний університет імені Тараса Шевченка \par}
{\parbox {5 cm}{Освітня програма \hfill}{\parbox {7 cm}{Прикладна математика\hfill}}\parbox {6 cm}{\hfill Семестр 1}}\par
{\parbox {5 cm} {Навчальний предмет \hfill}\parbox {7 cm}{Програмування \hfill}\parbox {6cm}{\hfill}\par}\vspace{-15pt} \end{scriptsize}}{\center Екзаменаційний білет \No \textbf { 41 }\par}}
{%beginitem
1. Компільовані та інтерпретовані мови.%enditem
\par
%beginitem
2. Написати програму, що запитує у користувача значення дійсних параметрів $x$ та $y$, обчислює для них значення виразу 
$$\frac{\ln x - 37}{(\tg x + 0.8) (y + 64)}$$ 
та повiдомляє введенi данi та результат обчислення.
 Оцінюється правильність та якість коду.

%enditem
\par
%beginitem
3. Написати функцію, що для цілого числа знаходить кількість простих дільників, які входять у розклад числа в найбільшому степені. Наприклад, для числа $3^5{\cdot}7^9{\cdot}11{\cdot}23^{9}$ таких дільників два – $23$ та $7$, тому функція має повернути $2$.  Використання рядків та структур даних заборонено. Якість та швидкодія коду також оцінюються. 

%enditem
}
{\smallskip\smallskip \begin {scriptsize} Затверджено на засіданні кафедри теоретичної кібернетики, протокол \No 1  від   06.09.2023 р.\par\smallskip\smallskip\smallskip\smallskip
{\parbox {6.5 cm} {Зав. кафедри \dotfill Ю.В.~Крак}}\hspace {0.7cm} {\hbox to 6.5 cm { Екзаменатор \dotfill Т.О.~Карнаух}} \end{scriptsize}}
%end
\vspace{0.9cm}
%begin
{{\begin {scriptsize}\par
{ \center  Київський національний університет імені Тараса Шевченка \par}
{\parbox {5 cm}{Освітня програма \hfill}{\parbox {7 cm}{Прикладна математика\hfill}}\parbox {6 cm}{\hfill Семестр 1}}\par
{\parbox {5 cm} {Навчальний предмет \hfill}\parbox {7 cm}{Програмування \hfill}\parbox {6cm}{\hfill}\par}\vspace{-15pt} \end{scriptsize}}{\center Екзаменаційний білет \No \textbf { 42 }\par}}
{%beginitem
1. Вбудовані числові типи мови Python.%enditem
\par
%beginitem
2. Написати програму, що запитує у користувача значення дійсних параметрів $x$ та $y$, обчислює для них значення виразу 
$$\frac{\log_{10} x + 68}{(\arccos x - 0.3) (y - 34)}$$ 
та повiдомляє введенi данi та результат обчислення.
 Оцінюється правильність та якість коду.

%enditem
\par
%beginitem
3. Написати функцію, що для інтервалу $[a, b)$ знаходить число з найбільшою кількістю різних простих дільників. Якщо їх кілька, то цікавить найменше.  Використання рядків та структур даних заборонено. Якість та швидкодія коду також оцінюються. 

%enditem
}
{\smallskip\smallskip \begin {scriptsize} Затверджено на засіданні кафедри теоретичної кібернетики, протокол \No 1  від   06.09.2023 р.\par\smallskip\smallskip\smallskip\smallskip
{\parbox {6.5 cm} {Зав. кафедри \dotfill Ю.В.~Крак}}\hspace {0.7cm} {\hbox to 6.5 cm { Екзаменатор \dotfill Т.О.~Карнаух}} \end{scriptsize}}
%end
\newpage
%begin
{{\begin {scriptsize}\par
{ \center  Київський національний університет імені Тараса Шевченка \par}
{\parbox {5 cm}{Освітня програма \hfill}{\parbox {7 cm}{Прикладна математика\hfill}}\parbox {6 cm}{\hfill Семестр 1}}\par
{\parbox {5 cm} {Навчальний предмет \hfill}\parbox {7 cm}{Програмування \hfill}\parbox {6cm}{\hfill}\par}\vspace{-15pt} \end{scriptsize}}{\center Екзаменаційний білет \No \textbf { 43 }\par}}
{%beginitem
1. Рекурсивні означення, рекурсивні функції.%enditem
\par
%beginitem
2. Написати програму, що запитує у користувача значення дійсних параметрів $x$ та $y$, обчислює для них значення виразу 
$$\frac{\ x - 54}{(\arccos x + 0.9) (y - 77)}$$ 
та повiдомляє введенi данi та результат обчислення.
 Оцінюється правильність та якість коду.

%enditem
\par
%beginitem
3. Написати функцію, що для цілого числа знаходить простий дільник, який входить у розклад числа в найбільшому степені. Якщо таких кілька, то повернути найменший. Наприклад, для числа $2^{9}{\cdot}3^5{\cdot}7^9{\cdot}11$ функція має повернути $2$.  Використання рядків та структур даних заборонено. Якість та швидкодія коду також оцінюються. 

%enditem
}
{\smallskip\smallskip \begin {scriptsize} Затверджено на засіданні кафедри теоретичної кібернетики, протокол \No 1  від   06.09.2023 р.\par\smallskip\smallskip\smallskip\smallskip
{\parbox {6.5 cm} {Зав. кафедри \dotfill Ю.В.~Крак}}\hspace {0.7cm} {\hbox to 6.5 cm { Екзаменатор \dotfill Т.О.~Карнаух}} \end{scriptsize}}
%end
\vspace{0.9cm}
%begin
{{\begin {scriptsize}\par
{ \center  Київський національний університет імені Тараса Шевченка \par}
{\parbox {5 cm}{Освітня програма \hfill}{\parbox {7 cm}{Прикладна математика\hfill}}\parbox {6 cm}{\hfill Семестр 1}}\par
{\parbox {5 cm} {Навчальний предмет \hfill}\parbox {7 cm}{Програмування \hfill}\parbox {6cm}{\hfill}\par}\vspace{-15pt} \end{scriptsize}}{\center Екзаменаційний білет \No \textbf { 44 }\par}}
{%beginitem
1. Об’єктна система Python. Поняття незмінюваного об'єкту.%enditem
\par
%beginitem
2. Написати програму, що запитує у користувача значення дійсних параметрів $x$ та $y$, обчислює для них значення виразу 
$$\frac{\ln x + 84}{(\cos x - 0.4) (y + 30)}$$ 
та повiдомляє введенi данi та результат обчислення.
 Оцінюється правильність та якість коду.

%enditem
\par
%beginitem
3. Написати функцію, яка для цілого числа знаходить суму його простих дільників, у десятковому запису яких нема цифри 3.  Використання рядків та структур даних заборонено. Якість та швидкодія коду також оцінюються. 

%enditem
}
{\smallskip\smallskip \begin {scriptsize} Затверджено на засіданні кафедри теоретичної кібернетики, протокол \No 1  від   06.09.2023 р.\par\smallskip\smallskip\smallskip\smallskip
{\parbox {6.5 cm} {Зав. кафедри \dotfill Ю.В.~Крак}}\hspace {0.7cm} {\hbox to 6.5 cm { Екзаменатор \dotfill Т.О.~Карнаух}} \end{scriptsize}}
%end
\vspace{0.9cm}
%begin
{{\begin {scriptsize}\par
{ \center  Київський національний університет імені Тараса Шевченка \par}
{\parbox {5 cm}{Освітня програма \hfill}{\parbox {7 cm}{Прикладна математика\hfill}}\parbox {6 cm}{\hfill Семестр 1}}\par
{\parbox {5 cm} {Навчальний предмет \hfill}\parbox {7 cm}{Програмування \hfill}\parbox {6cm}{\hfill}\par}\vspace{-15pt} \end{scriptsize}}{\center Екзаменаційний білет \No \textbf { 45 }\par}}
{%beginitem
1. Змінні в Python. Типізація. Операції над змінними. Присвоювання та його види у Python.%enditem
\par
%beginitem
2. Написати програму, що запитує у користувача значення дійсних параметрів $x$ та $y$, обчислює для них значення виразу 
$$\frac{\log_{10} x + 33}{(\tg x - 0.6) (y + 12)}$$ 
та повiдомляє введенi данi та результат обчислення.
 Оцінюється правильність та якість коду.

%enditem
\par
%beginitem
3. Написати функцію, що для інтервалу $[a, b)$ знаходить знаходить число з найбільшим добутком простих дільників. Якщо таких чисел кілька, повертає найменше.  Використання рядків та структур даних заборонено. Якість та швидкодія коду також оцінюються. 

%enditem
}
{\smallskip\smallskip \begin {scriptsize} Затверджено на засіданні кафедри теоретичної кібернетики, протокол \No 1  від   06.09.2023 р.\par\smallskip\smallskip\smallskip\smallskip
{\parbox {6.5 cm} {Зав. кафедри \dotfill Ю.В.~Крак}}\hspace {0.7cm} {\hbox to 6.5 cm { Екзаменатор \dotfill Т.О.~Карнаух}} \end{scriptsize}}
%end
\newpage
%begin
{{\begin {scriptsize}\par
{ \center  Київський національний університет імені Тараса Шевченка \par}
{\parbox {5 cm}{Освітня програма \hfill}{\parbox {7 cm}{Прикладна математика\hfill}}\parbox {6 cm}{\hfill Семестр 1}}\par
{\parbox {5 cm} {Навчальний предмет \hfill}\parbox {7 cm}{Програмування \hfill}\parbox {6cm}{\hfill}\par}\vspace{-15pt} \end{scriptsize}}{\center Екзаменаційний білет \No \textbf { 46 }\par}}
{%beginitem
1. Оператори порівняння (у широкому розумінні мови Python) та булеві оператори.%enditem
\par
%beginitem
2. Написати програму, що запитує у користувача значення дійсних параметрів $x$ та $y$, обчислює для них значення виразу 
$$\frac{\pi x - 27}{(\ctg x + 0.5) (y - 74)}$$ 
та повiдомляє введенi данi та результат обчислення.
 Оцінюється правильність та якість коду.

%enditem
\par
%beginitem
3. Написати функцію, що для інтервалу $[a, b)$ знаходить найближчу пару чисел з цього інтервалу, що мають парну кількість різних простих дільників. Якщо таких пар кілька, то цікавить пара з найбільшою сумою.  Використання рядків та структур даних заборонено. Якість та швидкодія коду також оцінюються. 

%enditem
}
{\smallskip\smallskip \begin {scriptsize} Затверджено на засіданні кафедри теоретичної кібернетики, протокол \No 1  від   06.09.2023 р.\par\smallskip\smallskip\smallskip\smallskip
{\parbox {6.5 cm} {Зав. кафедри \dotfill Ю.В.~Крак}}\hspace {0.7cm} {\hbox to 6.5 cm { Екзаменатор \dotfill Т.О.~Карнаух}} \end{scriptsize}}
%end
\vspace{0.9cm}
%begin
{{\begin {scriptsize}\par
{ \center  Київський національний університет імені Тараса Шевченка \par}
{\parbox {5 cm}{Освітня програма \hfill}{\parbox {7 cm}{Прикладна математика\hfill}}\parbox {6 cm}{\hfill Семестр 1}}\par
{\parbox {5 cm} {Навчальний предмет \hfill}\parbox {7 cm}{Програмування \hfill}\parbox {6cm}{\hfill}\par}\vspace{-15pt} \end{scriptsize}}{\center Екзаменаційний білет \No \textbf { 47 }\par}}
{%beginitem
1. Поняття часової та просторової складності обчислень. Складність у середньому. Асимптотичні оцінки складності.%enditem
\par
%beginitem
2. Написати програму, що запитує у користувача значення дійсних параметрів $x$ та $y$, обчислює для них значення виразу 
$$\frac{\sqrt x + 28}{(\sin x - 0.1) (y - 20)}$$ 
та повiдомляє введенi данi та результат обчислення.
 Оцінюється правильність та якість коду.

%enditem
\par
%beginitem
3. Написати функцію, що для цілого числа знаходить кількість його простих дільників, що входять у його розклад рівно у другому степеню. Наприклад, для числа $3^2{\cdot}7^9{\cdot}11{\cdot}23^{2}$ таких дільників два – $3$ та $23$, тому функція має повернути $2$.  Використання рядків та структур даних заборонено. Якість та швидкодія коду також оцінюються. 

%enditem
}
{\smallskip\smallskip \begin {scriptsize} Затверджено на засіданні кафедри теоретичної кібернетики, протокол \No 1  від   06.09.2023 р.\par\smallskip\smallskip\smallskip\smallskip
{\parbox {6.5 cm} {Зав. кафедри \dotfill Ю.В.~Крак}}\hspace {0.7cm} {\hbox to 6.5 cm { Екзаменатор \dotfill Т.О.~Карнаух}} \end{scriptsize}}
%end
\vspace{0.9cm}
%begin
{{\begin {scriptsize}\par
{ \center  Київський національний університет імені Тараса Шевченка \par}
{\parbox {5 cm}{Освітня програма \hfill}{\parbox {7 cm}{Прикладна математика\hfill}}\parbox {6 cm}{\hfill Семестр 1}}\par
{\parbox {5 cm} {Навчальний предмет \hfill}\parbox {7 cm}{Програмування \hfill}\parbox {6cm}{\hfill}\par}\vspace{-15pt} \end{scriptsize}}{\center Екзаменаційний білет \No \textbf { 48 }\par}}
{%beginitem
1. Емпіричний підхід до визначення складності обчислень.%enditem
\par
%beginitem
2. Написати програму, що запитує у користувача значення дійсних параметрів $x$ та $y$, обчислює для них значення виразу 
$$\frac{\log_{2} x - 13}{(\arctg x + 0.2) (y + 73)}$$ 
та повiдомляє введенi данi та результат обчислення.
 Оцінюється правильність та якість коду.

%enditem
\par
%beginitem
3. Написати функцію, що для інтервалу $[a, b)$ знаходить число з найбільшою кількістю різних простих дільників. Якщо їх кілька, то цікавить найбільше.  Використання рядків та структур даних заборонено. Якість та швидкодія коду також оцінюються. 

%enditem
}
{\smallskip\smallskip \begin {scriptsize} Затверджено на засіданні кафедри теоретичної кібернетики, протокол \No 1  від   06.09.2023 р.\par\smallskip\smallskip\smallskip\smallskip
{\parbox {6.5 cm} {Зав. кафедри \dotfill Ю.В.~Крак}}\hspace {0.7cm} {\hbox to 6.5 cm { Екзаменатор \dotfill Т.О.~Карнаух}} \end{scriptsize}}
%end
\newpage
%begin
{{\begin {scriptsize}\par
{ \center  Київський національний університет імені Тараса Шевченка \par}
{\parbox {5 cm}{Освітня програма \hfill}{\parbox {7 cm}{Прикладна математика\hfill}}\parbox {6 cm}{\hfill Семестр 1}}\par
{\parbox {5 cm} {Навчальний предмет \hfill}\parbox {7 cm}{Програмування \hfill}\parbox {6cm}{\hfill}\par}\vspace{-15pt} \end{scriptsize}}{\center Екзаменаційний білет \No \textbf { 49 }\par}}
{%beginitem
1. Глибина рекурсії та загальна кількість рекурсивних викликів, їх вплив на розміри програмного стеку та на час виконання програми.%enditem
\par
%beginitem
2. Написати програму, що запитує у користувача значення дійсних параметрів $x$ та $y$, обчислює для них значення виразу 
$$\frac{\sqrt x - 44}{(\arcsin x - 0.9) (y + 11)}$$ 
та повiдомляє введенi данi та результат обчислення.
 Оцінюється правильність та якість коду.

%enditem
\par
%beginitem
3. Написати функцію, що для цілого числа знаходить простий дільник, який входить у розклад числа в найбільшому степені. Якщо таких кілька, то повернути найбільший. Наприклад, для числа $2^{9}{\cdot}3^5{\cdot}7^9{\cdot}11$ функція має повернути $7$.  Використання рядків та структур даних заборонено. Якість та швидкодія коду також оцінюються. 

%enditem
}
{\smallskip\smallskip \begin {scriptsize} Затверджено на засіданні кафедри теоретичної кібернетики, протокол \No 1  від   06.09.2023 р.\par\smallskip\smallskip\smallskip\smallskip
{\parbox {6.5 cm} {Зав. кафедри \dotfill Ю.В.~Крак}}\hspace {0.7cm} {\hbox to 6.5 cm { Екзаменатор \dotfill Т.О.~Карнаух}} \end{scriptsize}}
%end
\vspace{0.9cm}
%begin
{{\begin {scriptsize}\par
{ \center  Київський національний університет імені Тараса Шевченка \par}
{\parbox {5 cm}{Освітня програма \hfill}{\parbox {7 cm}{Прикладна математика\hfill}}\parbox {6 cm}{\hfill Семестр 1}}\par
{\parbox {5 cm} {Навчальний предмет \hfill}\parbox {7 cm}{Програмування \hfill}\parbox {6cm}{\hfill}\par}\vspace{-15pt} \end{scriptsize}}{\center Екзаменаційний білет \No \textbf { 50 }\par}}
{%beginitem
1. Винятки та їх обробка.%enditem
\par
%beginitem
2. Написати програму, що запитує у користувача значення дійсних параметрів $x$ та $y$, обчислює для них значення виразу 
$$\frac{\ln x + 21}{(\arccos x + 0.7) (y - 95)}$$ 
та повiдомляє введенi данi та результат обчислення.
 Оцінюється правильність та якість коду.

%enditem
\par
%beginitem
3. Написати функцію, що для інтервалу $[a, b)$ знаходить найближчу пару чисел з цього інтервалу, що мають непарну кількість різних простих дільників. Якщо таких пар кілька, то цікавить пара з найменшим добутком.  Використання рядків та структур даних заборонено. Якість та швидкодія коду також оцінюються. 

%enditem
}
{\smallskip\smallskip \begin {scriptsize} Затверджено на засіданні кафедри теоретичної кібернетики, протокол \No 1  від   06.09.2023 р.\par\smallskip\smallskip\smallskip\smallskip
{\parbox {6.5 cm} {Зав. кафедри \dotfill Ю.В.~Крак}}\hspace {0.7cm} {\hbox to 6.5 cm { Екзаменатор \dotfill Т.О.~Карнаух}} \end{scriptsize}}
%end
\vspace{0.9cm}
%begin
{{\begin {scriptsize}\par
{ \center  Київський національний університет імені Тараса Шевченка \par}
{\parbox {5 cm}{Освітня програма \hfill}{\parbox {7 cm}{Прикладна математика\hfill}}\parbox {6 cm}{\hfill Семестр 1}}\par
{\parbox {5 cm} {Навчальний предмет \hfill}\parbox {7 cm}{Програмування \hfill}\parbox {6cm}{\hfill}\par}\vspace{-15pt} \end{scriptsize}}{\center Екзаменаційний білет \No \textbf { 51 }\par}}
{%beginitem
1. Лексеми мови Python.%enditem
\par
%beginitem
2. Написати програму, що запитує у користувача значення дійсних параметрів $x$ та $y$, обчислює для них значення виразу 
$$\frac{\log_{2} x + 64}{(\cos x + 0.4) (y - 45)}$$ 
та повiдомляє введенi данi та результат обчислення.
 Оцінюється правильність та якість коду.

%enditem
\par
%beginitem
3. Написати функцію, що для інтервалу $[a, b)$ знаходить найближчу пару чисел з цього інтервалу, що мають парну кількість різних простих дільників. Якщо таких пар кілька, то цікавить пара з найменшим добутком.  Використання рядків та структур даних заборонено. Якість та швидкодія коду також оцінюються. 

%enditem
}
{\smallskip\smallskip \begin {scriptsize} Затверджено на засіданні кафедри теоретичної кібернетики, протокол \No 1  від   06.09.2023 р.\par\smallskip\smallskip\smallskip\smallskip
{\parbox {6.5 cm} {Зав. кафедри \dotfill Ю.В.~Крак}}\hspace {0.7cm} {\hbox to 6.5 cm { Екзаменатор \dotfill Т.О.~Карнаух}} \end{scriptsize}}
%end
\newpage
%begin
{{\begin {scriptsize}\par
{ \center  Київський національний університет імені Тараса Шевченка \par}
{\parbox {5 cm}{Освітня програма \hfill}{\parbox {7 cm}{Прикладна математика\hfill}}\parbox {6 cm}{\hfill Семестр 1}}\par
{\parbox {5 cm} {Навчальний предмет \hfill}\parbox {7 cm}{Програмування \hfill}\parbox {6cm}{\hfill}\par}\vspace{-15pt} \end{scriptsize}}{\center Екзаменаційний білет \No \textbf { 52 }\par}}
{%beginitem
1. Глобальні, локальні та вільні змінні.%enditem
\par
%beginitem
2. Написати програму, що запитує у користувача значення дійсних параметрів $x$ та $y$, обчислює для них значення виразу 
$$\frac{\pi x - 71}{(\arctg x - 0.6) (y + 84)}$$ 
та повiдомляє введенi данi та результат обчислення.
 Оцінюється правильність та якість коду.

%enditem
\par
%beginitem
3. Написати функцію, що для інтервалу $[a, b)$ знаходить число з найбільшою кількістю різних простих дільників, що містять цифру 7 у своєму десятковому запису. Якщо їх кілька, то цікавить найменше.  Використання рядків та структур даних заборонено. Якість та швидкодія коду також оцінюються. 

%enditem
}
{\smallskip\smallskip \begin {scriptsize} Затверджено на засіданні кафедри теоретичної кібернетики, протокол \No 1  від   06.09.2023 р.\par\smallskip\smallskip\smallskip\smallskip
{\parbox {6.5 cm} {Зав. кафедри \dotfill Ю.В.~Крак}}\hspace {0.7cm} {\hbox to 6.5 cm { Екзаменатор \dotfill Т.О.~Карнаух}} \end{scriptsize}}
%end
\vspace{0.9cm}
%begin
{{\begin {scriptsize}\par
{ \center  Київський національний університет імені Тараса Шевченка \par}
{\parbox {5 cm}{Освітня програма \hfill}{\parbox {7 cm}{Прикладна математика\hfill}}\parbox {6 cm}{\hfill Семестр 1}}\par
{\parbox {5 cm} {Навчальний предмет \hfill}\parbox {7 cm}{Програмування \hfill}\parbox {6cm}{\hfill}\par}\vspace{-15pt} \end{scriptsize}}{\center Екзаменаційний білет \No \textbf { 53 }\par}}
{%beginitem
1. Математичний підхід до визначення складності обчислень.%enditem
\par
%beginitem
2. Написати програму, що запитує у користувача значення дійсних параметрів $x$ та $y$, обчислює для них значення виразу 
$$\frac{\ x - 58}{(\tg x - 0.3) (y - 40)}$$ 
та повiдомляє введенi данi та результат обчислення.
 Оцінюється правильність та якість коду.

%enditem
\par
%beginitem
3. Написати функцію, що для інтервалу $[a, b)$ знаходить число з найбільшою кількістю різних простих дільників, що містять цифру 7 у своєму десятковому запису. Якщо їх кілька, то цікавить найменше.  Використання рядків та структур даних заборонено. Якість та швидкодія коду також оцінюються. 

%enditem
}
{\smallskip\smallskip \begin {scriptsize} Затверджено на засіданні кафедри теоретичної кібернетики, протокол \No 1  від   06.09.2023 р.\par\smallskip\smallskip\smallskip\smallskip
{\parbox {6.5 cm} {Зав. кафедри \dotfill Ю.В.~Крак}}\hspace {0.7cm} {\hbox to 6.5 cm { Екзаменатор \dotfill Т.О.~Карнаух}} \end{scriptsize}}
%end
\vspace{0.9cm}
%begin
{{\begin {scriptsize}\par
{ \center  Київський національний університет імені Тараса Шевченка \par}
{\parbox {5 cm}{Освітня програма \hfill}{\parbox {7 cm}{Прикладна математика\hfill}}\parbox {6 cm}{\hfill Семестр 1}}\par
{\parbox {5 cm} {Навчальний предмет \hfill}\parbox {7 cm}{Програмування \hfill}\parbox {6cm}{\hfill}\par}\vspace{-15pt} \end{scriptsize}}{\center Екзаменаційний білет \No \textbf { 54 }\par}}
{%beginitem
1. Цикл while, умова продовження циклу, вигляд і семантика інструкції циклу while.%enditem
\par
%beginitem
2. Написати програму, що запитує у користувача значення дійсних параметрів $x$ та $y$, обчислює для них значення виразу 
$$\frac{\log_{10} x + 76}{(\arcsin x + 0.8) (y + 86)}$$ 
та повiдомляє введенi данi та результат обчислення.
 Оцінюється правильність та якість коду.

%enditem
\par
%beginitem
3. Написати функцію, що для інтервалу $[a, b)$ знаходить найближчу пару чисел з цього інтервалу, що мають непарну кількість різних простих дільників. Якщо таких пар кілька, то цікавить пара з найбільшою сумою.  Використання рядків та структур даних заборонено. Якість та швидкодія коду також оцінюються. 

%enditem
}
{\smallskip\smallskip \begin {scriptsize} Затверджено на засіданні кафедри теоретичної кібернетики, протокол \No 1  від   06.09.2023 р.\par\smallskip\smallskip\smallskip\smallskip
{\parbox {6.5 cm} {Зав. кафедри \dotfill Ю.В.~Крак}}\hspace {0.7cm} {\hbox to 6.5 cm { Екзаменатор \dotfill Т.О.~Карнаух}} \end{scriptsize}}
%end
\newpage
\end{document}
